% \clearpage
\section{Conclusion}
This paper quantifies the buildup of systemic risks in the Armenian economy, differentiates between those risks and considers their effects on future output growth. It also measures the impact of macroprudential policy on systemic risks and uses two approaches to derive the macroprudential stance. By employing both panel and localized models, the study ensures robust results, accounting for the nonlinear dynamics of systemic risks and policy measures. The framework integrates financial cycles, financial conditions, liquidity and capital risks, macroeconomic factors, and their interactions with macroprudential policy, while also considering the associated costs of the policy.

The findings reveal that the lag between systemic risk buildup and materialization is shorter in Armenia compared to other countries, with financial vulnerabilities following a cyclical trend and financial stress emerging more abruptly. Macroprudential policy affects the resilience of the system and is most effective during periods of low financial stress. The most significant risk reductions occur when tightening policies are applied during phases of rapid credit growth, however, this comes with larger short-term costs. Historically, the impact of the policy on financial vulnerabilities is much stronger than on financial stress.

As a small, open economy, Armenia’s systemic risks are more sensitive to external shocks, but these externalities affect the entire output growth distribution rather uniformly. This suggests that while external shocks are a concern, their direct impact on systemic risks remains small.

Using the distance-to-tail approach, we estimate that a policy tightening threshold of around 7\% is appropriate for Armenia, whereas cross-country comparisons suggest a range of 5-6\%.\footnote{\autoref{sec:dtt}} Additionally, when considering the welfare function approach, we estimate that the historical relative risk aversion in Armenia is around 3.57\footnote{\ref{appx:w_calibration}} and offer insights on how to refine policy decisions. This framework aims to improve the decision-making process of the policymaker by defining clear guidelines for preemptive action.